\subsection{关键实现方法}

\subsubsection{多控制相位翻转}

Oracle和Diffusion都需要对\(|11\cdots 1\rangle\)执行相位翻转。程序将这个操作封装为\texttt{\_apply\_multi\_controlled\_z}。当量子比特数为1时直接使用\(Z\)门；当量子比特数大于1时，程序用\(H-\mathrm{MCX}-H\)把多控制\(X\)门转化为多控制\(Z\)门：

\begin{lstlisting}[language=Python]
def _apply_multi_controlled_z(qc, qubits):
    if len(qubits) == 1:
        qc.z(qubits[0])
        return

    controls = qubits[:-1]
    target = qubits[-1]
    qc.h(target)
    qc.mcx(controls, target)
    qc.h(target)
\end{lstlisting}

以3量子比特为例，该函数会对\(|111\rangle\)态乘以\(-1\)，其余计算基态保持不变。

\subsubsection{目标态Oracle构造}

项目允许用户选择任意目标态。Oracle的目标是只对该目标态执行相位翻转。实现时，程序先把目标态临时映射到\(|11\cdots 1\rangle\)，再调用多控制相位翻转，最后恢复原来的计算基顺序。

例如目标态为\(|101\rangle\)时，中间的第2位为0，程序会先对对应量子比特施加\(X\)门，把\(|101\rangle\)映射到\(|111\rangle\)。完成多控制\(Z\)门后，再施加一次\(X\)门还原：

\begin{lstlisting}[language=Python]
def apply_oracle(qc, target, add_barrier=True):
    n = len(target)

    for qubit, bit in enumerate(target):
        if bit == "0":
            qc.x(qubit)

    _apply_multi_controlled_z(qc, list(range(n)))

    for qubit, bit in enumerate(target):
        if bit == "0":
            qc.x(qubit)

    if add_barrier:
        qc.barrier()
\end{lstlisting}

因此，对于任意用户输入的目标比特串，Oracle都会实现
\[
O|\omega\rangle=-|\omega\rangle,\qquad O|x\rangle=|x\rangle\quad(x\ne \omega).
\]

\subsubsection{Diffusion算子构造}

Diffusion算子采用标准的\(H-X-\mathrm{MCZ}-X-H\)结构。程序先用Hadamard门和\(X\)门把关于均匀叠加态的反演转换为对\(|11\cdots 1\rangle\)的相位翻转，再恢复基底：

\begin{lstlisting}[language=Python]
def apply_diffusion(qc, n, add_barrier=True):
    for qubit in range(n):
        qc.h(qubit)
    for qubit in range(n):
        qc.x(qubit)

    _apply_multi_controlled_z(qc, list(range(n)))

    for qubit in range(n):
        qc.x(qubit)
    for qubit in range(n):
        qc.h(qubit)

    qc.global_phase += np.pi

    if add_barrier:
        qc.barrier()
\end{lstlisting}

其中，\texttt{qc.global\_phase += np.pi}用于校正全局相位。经过该处理后，保存的statevector与
\[
    a_i'=2\bar{a}-a_i
\]
一致，Diffusion后的概率幅可以直接解释为关于平均概率幅的反演。

\subsubsection{完整Grover电路构造}

完整电路由Hadamard初始化和若干轮Grover迭代组成。程序在\texttt{build\_grover\_circuit}中先创建量子电路，再对全部量子比特施加Hadamard门，随后按迭代次数重复调用Oracle和Diffusion：

\begin{lstlisting}[language=Python]
def build_grover_circuit(n, target, iterations,
                         measure=False, add_barriers=True):
    validate_problem(n, target, iterations)

    if measure:
        qc = QuantumCircuit(n, n)
    else:
        qc = QuantumCircuit(n)

    qc.h(range(n))
    if add_barriers:
        qc.barrier()

    for _ in range(iterations):
        apply_oracle(qc, target, add_barrier=add_barriers)
        apply_diffusion(qc, n, add_barrier=add_barriers)

    if measure:
        for qubit in range(n):
            qc.measure(qubit, n - 1 - qubit)

    return qc
\end{lstlisting}

测量模式下，程序把第\(i\)个量子比特测量到第\(n-1-i\)个经典比特，使最终输出的比特串顺序与用户输入的目标态顺序一致。

\subsubsection{中间态保存}

为了得到每一步后的概率幅数据，程序没有只保存最终电路，而是在\texttt{build\_grover\_history}中逐步构造电路。每完成Hadamard初始化、Oracle或Diffusion操作，就复制当前电路，并用\texttt{Statevector.from\_instruction(qc)}计算当前量子态：

\begin{lstlisting}[language=Python]
def build_grover_history(n, target, iterations):
    validate_problem(n, target, iterations)

    qc = QuantumCircuit(n)
    history = []

    qc.h(range(n))
    qc.barrier()
    history.append(GroverStep(
        index=0,
        kind="initial",
        iteration=0,
        circuit=qc.copy(),
        statevector=Statevector.from_instruction(qc),
    ))

    step_index = 1
    for iteration in range(1, iterations + 1):
        apply_oracle(qc, target, add_barrier=True)
        history.append(GroverStep(
            index=step_index,
            kind="oracle",
            iteration=iteration,
            circuit=qc.copy(),
            statevector=Statevector.from_instruction(qc),
        ))
        step_index += 1

        apply_diffusion(qc, n, add_barrier=True)
        history.append(GroverStep(
            index=step_index,
            kind="diffusion",
            iteration=iteration,
            circuit=qc.copy(),
            statevector=Statevector.from_instruction(qc),
        ))
        step_index += 1

    return history
\end{lstlisting}

该历史序列直接作为后续概率幅图、概率图和3D步进图的数据来源。
